\chapter{Replacement cost regression}
\label{app:cost}

Multiple linear regression relates the replacement cost of six Dutch quay wall
projects (Table~\ref{tab:dutch-quay-walls}) to their length and to a heritage
and a disruption score, to test whether a cost that depends on the site would
be better than the cost per metre of Section~\ref{sec:tradeoff}. Lengths, works
and costs come from public sources; the disruption and heritage scores, from 1
(low) to 5 (high), were judged by the team from documented site
characteristics and are treated as equally spaced numbers.

\begin{table}[htbp]
\centering
\tiny
\caption{Characteristics and approximate replacement costs of regression quay-wall projects.}
\label{tab:dutch-quay-walls}
\setlength{\tabcolsep}{3pt}
\renewcommand{\arraystretch}{1.2}
\begin{tabular}{@{} l l r l c l c r l @{}}
\toprule
Quay wall & Replacement works & Len. & Disruption factors & Disp. & Heritage factors & Her. & Cost & References \\
\midrule
Beestenmarkt, & Full replacement of masonry & 12 & Very restricted access and & 5 & Protected historic townscape, & 4 & 0.52 & \cite{waterland2025,} \\
Monnickendam & quay; reuse of bricks. & & vehicle dimensions; pontoons. & & adjacent to historic buildings. & & & \cite{kdekker2025} \\
\addlinespace
Maaskade, & Reconstruction and raising of & 500 & Restricted parking, waterfront & 4 & Late 19th-c. quay, historic & 4 & 8.03 & \cite{maaskadeaward2018,} \\
Rotterdam & existing quay and drainage. & & works and berth loss. & & streetscape and cobblestones.  \\
\addlinespace
Haven vak M2, & Complete quay-wall and & 210 & Works in active historic harbour; & 5 & Protected townscape, demolition & 4 & 7.89 & \cite{maassluis2026,} \\
Maassluis & foundation replacement. & & public/waterway interfaces. & & permit in protected area. & & & \\
\addlinespace
Tuingracht, & Planned full replacement of & 330 & Traffic loading, parked vehicles & 4 & Nationally designated protected & 5 & 3.04 & \cite{alkmaarMIP,} \\
De Rijp & masonry wall on south bank. & & and constrained centre access. & & village centre; historic quay. & & & \cite{tuingrachtTender} \\
\addlinespace
Houtzagerssingel, & Demolition and replacement; & $\sim$370 & Utility/sewer diversions, tree & 4 & Established urban canal setting; & 3 & 6.26 & \cite{houtzagersTender,} \\
Den Haag & sheet piles, anchors, paving. & & works, pavement reinstatement. & & reinstatement of features. & & & \cite{houtzagersContractor} \\
\addlinespace
Westerkade West, & Construction of new retaining & 150 & Quayside construction and local & 3 & Historic quay dating from 1922 & 3 & 4.75 & \cite{westerkadeTender,} \\
Rotterdam & wall with steel sheet piles. & & paving/access changes. & & and waterfront character. & & & \cite{westerkadeContractor} \\
\bottomrule
\end{tabular}
\end{table}

The model explains 63.7\,\% of the variation in cost ($R^2 = 0.637$), but with
six projects and three predictors only two degrees of freedom remain for the
residuals: the adjusted $R^2$ is 0.094, and no predictor is significant at
5\,\% (length $p = 0.218$, heritage $p = 0.429$, disruption $p = 0.606$).
Forward and backward selection by AIC keep only the length ($R^2 = 0.457$), and
for a wall of 55\,m its 95\,\% prediction interval runs from \textminus\euro5.4
to \euro10.8 million. Section~\ref{sec:tradeoff} therefore keeps the cost of
about \euro35k per metre from the programme plan
\parencite{utrecht_programmaplan}.
